% TEMPLATE-MILC-v3.tex - plantilla maestra UNICA de las guias MILC v3.
%
% La usan las tres familias de guias, cada una con su driver:
%   · Grados 6-11  -> scripts/build-guias-g11.py
%   · Bebras       -> scripts/build-guias-bebras.py
%   · Semillero    -> scripts/build-guias-semillero.py
%
% Antes habia tres copias casi identicas (~400 lineas iguales, 6 distintas):
% cada mejora habia que aplicarla tres veces, y en la practica no se hacia
% (el motor de recursos vivia solo en dos de ellas). Lo que varia entre
% familias esta parametrizado en 6 placeholders que cada driver llena
% (se nombran aqui SIN su sintaxis real: el builder reemplaza por texto plano,
% tambien dentro de los comentarios, y un valor multilinea rompe el preambulo):
%
%   HEADER_IZQ        encabezado izquierdo de pagina
%   HEADER_DER        encabezado derecho de pagina
%   PORTADA_ETIQUETA  rotulo sobre el numero grande (GRADO/MOMENTO/MODULO)
%   PORTADA_SERIE     linea de serie de la portada
%   PORTADA_ID        identificador de la guia en la portada
%   PORTADA_CIRCULO   el circulito de grado (°), o vacio si no aplica
%
% El resto de claves las documenta content/guias/_SCHEMA.md. El script valida
% que no queden placeholders sin reemplazar antes de escribir el archivo final.

\documentclass[11pt,letterpaper]{article}
\usepackage[margin=2.0cm]{geometry}
\usepackage[table]{xcolor}
\usepackage{fontspec}
\usepackage[spanish]{babel}
\usepackage{tikz}
% Librerías TikZ para los diagramas de las guías (bloque `recursos:`):
%  · babel  → babel en español vuelve activos caracteres como «>», y TikZ los usa
%             en sus opciones (>=stealth, ->). Sin esta librería, cualquier
%             diagrama con flechas revienta con «\language@active@arg> has an
%             extra }». Debe ir ANTES de que se dibuje nada.
%  · positioning        → sintaxis «below left=2pt and 3pt».
%  · decorations.pathreplacing → llaves ({brace}) para señalar rangos.
\usetikzlibrary{babel,positioning,decorations.pathreplacing}
\usepackage{graphicx}
% Circuitos: el trabajo de electrónica se hace hoy con micro:bit y sus sensores
% integrados (MakeCode, sin cableado), así que no se necesitan esquemáticos.
% Si algún día entran montajes con protoboard, basta descomentar esta línea:
% los diagramas .tex/.tikz declarados en `recursos:` ya se incrustan solos.
% \usepackage{circuitikz}
\usepackage[most]{tcolorbox}
\usepackage{tabularx}
\usepackage{array}
\usepackage{fancyhdr}
\usepackage{titlesec}
\usepackage[document]{ragged2e}  % bandera a la derecha en todo el cuerpo
\usepackage{enumitem}
\usepackage{microtype}

% Helvetica Neue no trae la flecha U+2192, asi que cada «→» escrito literalmente
% en el YAML se compilaba a NADA, en silencio (462 flechas en 61 guias: rutas del
% tipo «paleta → tipografia → lockup» salian rotas). Mapeamos el caracter a su
% version matematica, que si tiene glifo. Asi el autor puede seguir escribiendo
% «→» comodamente en el YAML.
\usepackage{newunicodechar}
\newunicodechar{→}{\ensuremath{\rightarrow}}
% Mismo problema con los simbolos de marca y vinetas: el «✓» de «una palomita ✓»
% salia como cuadrito vacio en el PDF de todas las guias que lo usaban.
\usepackage{pifont}
\newunicodechar{✓}{\ding{51}}
\newunicodechar{✗}{\ding{55}}
\newunicodechar{✔}{\ding{52}}
\newunicodechar{•}{\textbullet}
\newunicodechar{≥}{\ensuremath{\geq}}
\newunicodechar{≤}{\ensuremath{\leq}}

\setmainfont{Helvetica Neue}
\setsansfont{Helvetica Neue}
\setlength{\parindent}{0pt}
\setlength{\parskip}{6pt}
% Sin particion de palabras: en 6.º-7.º una palabra cortada por guion al final
% de linea («Comuni-cacion») frena la lectura mucho mas que una linea corta.
\hyphenpenalty=10000
\exhyphenpenalty=10000
\pretolerance=2000
\tolerance=3000
\setlength{\emergencystretch}{3em}
\linespread{1.10}
\renewcommand{\arraystretch}{1.25}
\setlist{leftmargin=5mm,itemsep=1mm,topsep=1mm}
\newcolumntype{Y}{>{\RaggedRight\arraybackslash}X}

\definecolor{milcMagenta}{HTML}{E6007E}
\definecolor{milcVino}{HTML}{5A0038}
\definecolor{milcTurquesa}{HTML}{0093A5}
\definecolor{milcVerde}{HTML}{58A923}
\definecolor{milcMostaza}{HTML}{E5B400}
\definecolor{milcOcre}{HTML}{B86600}
\definecolor{milcNegro}{HTML}{241816}
\definecolor{milcGris}{HTML}{F7F7F4}
\definecolor{milcLinea}{HTML}{E8E2DD}
\definecolor{milcGradePrimary}{HTML}{FF2D87}
\definecolor{milcGradeDark}{HTML}{9F1B56}
\definecolor{milcGradeSoft}{HTML}{FFE0EE}
\definecolor{milcGradeAccent}{HTML}{CFE7FF}
\definecolor{milcGradeText}{HTML}{FFFFFF}
\color{milcNegro}

\pagestyle{fancy}
\fancyhf{}
\lhead{\small Tecnología e Informática · Grado 10}
\rhead{\small Guía 10 · MILC v3}
\cfoot{\small\thepage}
\renewcommand{\headrulewidth}{0pt}

\newtcolorbox{titlebox}[1]{
  enhanced,colback=#1,colframe=#1,arc=7mm,boxrule=0pt,
  left=7mm,right=7mm,top=3mm,bottom=3mm,
  before skip=8pt,after skip=8pt,
  fontupper=\sffamily\bfseries\Large\color{white}
}

\newtcolorbox{softbox}[3]{
  enhanced,colback=#2,colframe=#1,borderline west={4pt}{0pt}{#1},
  arc=2mm,boxrule=.4pt,left=4mm,right=4mm,top=3mm,bottom=3mm,
  before skip=6pt,after skip=8pt,title={#3},coltitle=white,
  colbacktitle=#1,fonttitle=\sffamily\bfseries,fontupper=\RaggedRight,
  attach boxed title to top left={xshift=3mm,yshift=-2mm},
  boxed title style={arc=2mm, boxrule=0pt}
}

\newtcolorbox{drawbox}[2]{
  enhanced,colback=white,colframe=#1,arc=4mm,boxrule=1pt,
  left=5mm,right=5mm,top=4mm,bottom=4mm,before skip=8pt,after skip=8pt,
  title={#2},coltitle=white,colbacktitle=#1,fonttitle=\sffamily\bfseries,
  fontupper=\RaggedRight,
  attach boxed title to top left={xshift=4mm,yshift=-2mm},
  boxed title style={arc=3mm, boxrule=0pt}
}

\newtcolorbox{infoband}[1]{
  enhanced,colback=#1!8,colframe=#1!50,arc=2mm,boxrule=.5pt,
  left=4mm,right=4mm,top=2mm,bottom=2mm,before skip=4pt,after skip=4pt,
  fontupper=\small\RaggedRight
}

% Puente narrativo entre secciones (contrato editorial Conéctate).
% Da continuidad: "Acabas de X. Ahora vas a Y porque Z."
% Visualmente: banda izquierda en color del grado, texto en itálica pequeña.
\newtcolorbox{puentebox}{
  enhanced,colback=milcGradeSoft,colframe=milcGradeDark,
  borderline west={3pt}{0pt}{milcGradeDark},
  arc=0mm,boxrule=0pt,left=5mm,right=4mm,top=2.5mm,bottom=2.5mm,
  before skip=4pt,after skip=8pt,
  fontupper=\itshape\small\color{milcNegro}\RaggedRight
}
% La flecha va en modo matematico a proposito: Helvetica Neue NO tiene el glifo
% U+2192, asi que \textrightarrow se compilaba a NADA («Missing character: There
% is no → in font Helvetica Neue») y el puente salia sin flecha. En modo
% matematico la flecha viene de la fuente matematica y si se dibuja.
\newcommand{\puente}[1]{%
  \begin{puentebox}%
    \textcolor{milcGradeDark}{$\rightarrow$}\hspace{2mm}#1%
  \end{puentebox}%
}

\newcommand{\linead}[1]{\noindent\color{milcLinea}\rule{#1}{0.5pt}\color{milcNegro}}
\newcommand{\respuesta}[1]{\par\vspace{2mm}\foreach \n in {1,...,#1}{\noindent\color{milcLinea}\rule{\linewidth}{0.5pt}\par\vspace{5mm}}\color{milcNegro}}
\newcommand{\checkbox}{\textcolor{milcGradeDark}{\fbox{\phantom{x}}}\hspace{5pt}}

\newcommand{\phasecard}[4]{
\begin{tcolorbox}[enhanced,colback=#3,colframe=#2,arc=3mm,boxrule=.5pt,height=3.05cm,valign=center,halign=center,left=1.5mm,right=1.5mm,top=2mm,bottom=2mm]
{\sffamily\bfseries\fontsize{22}{24}\selectfont\textcolor{#2}{#1}}\par
{\sffamily\bfseries\small #4}
\end{tcolorbox}
}

% ── Piezas del rediseno 2026 (legibilidad 6.º-11.º) ───────────────────
% Banda de estacion: reemplaza el titlebox macizo. Titulo a la izquierda,
% dato practico (tiempo/modalidad) a la derecha, en una sola linea.
\newcommand{\milcestacion}[2]{%
\begin{tcolorbox}[enhanced,colback=#1,colframe=#1,arc=2mm,boxrule=0pt,
  left=5mm,right=5mm,top=2.6mm,bottom=2.6mm,before skip=11pt,after skip=7pt]
{\sffamily\bfseries\fontsize{15}{18}\selectfont\color{white}#2}
\end{tcolorbox}}

% Tarjeta de paso numerada: sustituye las tablas «Pilar / Aplicacion», que
% eran formato de consulta docente, no de estudiante. El numero va en un
% circulo sobre el borde para que la secuencia se lea de un vistazo.
\newcommand{\milcpaso}[3]{%
\begin{tcolorbox}[enhanced,colback=white,colframe=milcLinea,arc=1.5mm,boxrule=.9pt,
  left=14mm,right=4mm,top=3mm,bottom=3mm,before skip=4pt,after skip=4pt,
  fontupper=\RaggedRight,
  overlay={\node[anchor=north west,circle,fill=#2,text=white,inner sep=0pt,
    minimum size=7.4mm,font=\sffamily\bfseries\small]
    at ([xshift=3.6mm,yshift=-3.2mm]frame.north west) {#1};}]
#3
\end{tcolorbox}}

% Casilla marcable de tamano util con lapiz (la anterior era \fbox{\phantom{x}},
% de alto variable segun el texto que la seguia).
\newcommand{\milcmarca}{\framebox[3.8mm]{\rule{0pt}{3.4mm}}}

% Fila marcable de lista suelta (checks de la estacion 1).
\newcommand{\milccheckrow}[1]{%
\par\vspace{1.8mm}\noindent
\makebox[6.4mm][l]{\raisebox{-.55mm}{\textcolor{milcNegro}{\milcmarca}}}%
\parbox[t]{\dimexpr\linewidth-6.4mm\relax}{\RaggedRight #1}\par}

% Celda marcable dentro de la rubrica (tabla de autoevaluacion). Misma
% sangria francesa, pero pensada para vivir dentro de una columna X.
\newcommand{\milcopcion}[1]{%
\makebox[5.2mm][l]{\raisebox{-.55mm}{\textcolor{milcNegro}{\milcmarca}}}%
\parbox[t]{\dimexpr\linewidth-5.2mm\relax}{\RaggedRight #1}}

% ── Recursos visuales de la guía ──────────────────────────────────────
% Declarados en el bloque `recursos:` del YAML y emitidos por el builder.
% \guiaFigura   → imagen rasterizada o PDF (foto, carta celeste, montaje).
% \guiaDiagrama → fuente TikZ/circuitikz (.tex) incrustada como vector:
%                 es la vía para circuitos y esquemáticos editables.
% #1 = ruta relativa al .tex · #2 = pie (puede ir vacío)
\newcommand{\guiaFigura}[2]{%
\begin{center}
\includegraphics[width=0.88\linewidth,height=0.38\textheight,keepaspectratio]{#1}\\[1.5mm]
{\footnotesize\itshape\textcolor{milcNegro!75}{#2}}
\end{center}
}

\newcommand{\guiaDiagrama}[2]{%
\begin{center}
\input{#1}\\[1.5mm]
{\footnotesize\itshape\textcolor{milcNegro!75}{#2}}
\end{center}
}

\begin{document}
\thispagestyle{empty}

% ── Portada ───────────────────────────────────────────────────────────
% Antes: un rectangulo de color a pagina completa. Se veia bien en pantalla,
% pero en la fotocopiadora del colegio se comia el toner y salia gris sucio.
% Ahora la banda ocupa el tercio superior y el resto va en blanco.
\begin{tikzpicture}[remember picture,overlay]
  \path[fill=milcGradePrimary]
    (current page.north west) rectangle ([yshift=-10.6cm]current page.north east);

  \node[anchor=north west,text=milcGradeText] at ([xshift=2.0cm,yshift=-1.55cm]current page.north west)
    {\sffamily\bfseries\fontsize{12}{14}\selectfont GRADO};
  \node[anchor=north west,text=milcGradeText,opacity=.85] at ([xshift=2.0cm,yshift=-2.35cm]current page.north west)
    {\sffamily\bfseries\fontsize{8.5}{10}\selectfont SERIE GUÍAS · TECNOLOGÍA E INFORMÁTICA};
  \node[anchor=north east,text=milcGradeText,opacity=.85,align=right] at ([xshift=-2.0cm,yshift=-1.55cm]current page.north east)
    {\sffamily\bfseries\fontsize{9}{12}\selectfont I.E. Sor María Juliana\\Cartago, Valle del Cauca};

  \node[anchor=west,text=milcGradeText] (gradonum) at ([xshift=1.85cm,yshift=-7.1cm]current page.north west)
    {\sffamily\bfseries\fontsize{112}{112}\selectfont 10};
  % El circulito de grado (°) solo aplica a las guias de grado («11°»); en
  % Bebras y Semillero el numero grande es un momento/modulo, no un grado.
  % Se posiciona relativo al nodo `gradonum`, no en coordenadas absolutas.
  \draw[milcGradeText,line width=1.6pt]
    ([xshift=2.0mm,yshift=-4.5mm]gradonum.north east) circle (.15cm);

  \node[anchor=west,text=milcGradeText,text width=11.4cm,align=left] at ([xshift=1.5cm,yshift=0.5cm]gradonum.east)
    {
     {\sffamily\bfseries\fontsize{9}{11}\selectfont\MakeUppercase{Período 1 · Escritura abierta con IA}}\\[3mm]
     {\sffamily\bfseries\fontsize{21}{25}\selectfont\setlength{\baselineskip}{27pt}Sustentación +\\[4pt]feria del libro\\[4pt]escolar}\\[3.5mm]
     {\sffamily\bfseries\fontsize{15}{18}\selectfont Guía 10-10-TIC}
    };
\end{tikzpicture}

\vspace*{9.4cm}
\vfill

{\sffamily\bfseries\fontsize{8}{10}\selectfont\textcolor{milcGradeDark}{LO QUE TE LLEVAS HOY}}

\begin{tcolorbox}[enhanced,colback=white,colframe=milcNegro,arc=2mm,boxrule=1.6pt,
  left=5mm,right=5mm,top=4mm,bottom=4mm,before skip=3pt,after skip=12pt,
  fontupper=\RaggedRight\fontsize{11}{15.5}\selectfont]
Sustentación pública de 6 minutos de tu libro frente al grupo con (a) presentación del libro (concepción + proceso + uso de IA + decisiones críticas), (b) lectura en voz alta de 1-2 minutos de un fragmento elegido, (c) muestra del PDF o libro impreso, (d) sesión de preguntas y respuestas de 5 minutos + participación activa en feria del libro escolar (puesto del libro, lectura de fragmentos de compañeros, intercambio honesto de comentarios) + autoevaluación final con 2 fortalezas y 2 mejoras
\end{tcolorbox}

\begin{tcolorbox}[enhanced,colback=milcGris,colframe=milcGris,arc=2mm,boxrule=0pt,
  left=5mm,right=5mm,top=3.5mm,bottom=3.5mm,after skip=12pt]
{\sffamily\bfseries\fontsize{8}{10}\selectfont\textcolor{milcNegro!55}{ESTA GUÍA ES DE}}\\[3mm]
\begin{tabularx}{\linewidth}{X p{3.2cm} p{3.2cm}}
\linead{\linewidth} & \linead{3.0cm} & \linead{3.0cm}\\[-1mm]
{\fontsize{8.5}{10}\selectfont\textcolor{milcNegro!55}{Nombre completo}} &
{\fontsize{8.5}{10}\selectfont\textcolor{milcNegro!55}{Grupo}} &
{\fontsize{8.5}{10}\selectfont\textcolor{milcNegro!55}{Fecha}}\\
\end{tabularx}
\end{tcolorbox}

\vfill

\noindent\textcolor{milcLinea}{\rule{\linewidth}{1pt}}\\[2mm]
{\sffamily\bfseries\fontsize{9.5}{12}\selectfont Autor: Dr. Álvaro Cárdenas Orozco}\\[1mm]
{\sffamily\fontsize{8.5}{11}\selectfont\textcolor{milcNegro!55}{Metodología MILC · apertura ancestral · pensamiento computacional · triángulo Dussel–estoicismo–Floridi}}

\newpage

\section*{Apertura: Sustentación + feria del libro escolar}

\begin{softbox}{milcGradeDark}{milcGradeSoft}{Saber ancestral}
En los pueblos del Valle del Cauca y en los corregimientos cafeteros del Quindío, durante décadas existió una práctica cultural que reunía cada año al pueblo entero: \textbf{la feria del libro del pueblo}. No era feria de editorial grande: era encuentro local. Los autores barriales (poetas aficionados, cronistas del pueblo, maestros que habían escrito libros pedagógicos, párrocos con sermones recopilados) traían sus libros impresos en pequeño tiraje al parque principal. Ponían mesa con mantel sencillo, exhibían los ejemplares, se sentaban a esperar. Pero la feria no era solo exhibición: era \textbf{intercambio activo}. Los autores leían fragmentos en voz alta cuando llegaba alguien interesado. Los lectores preguntaban: \emph{``¿de qué trata?''}, \emph{``¿por qué lo escribió?''}, \emph{``¿está pasando lo que cuenta?''}. Los libros se vendían a precios bajos, se prestaban, se intercambiaban entre autores. Al cerrar el día, cada autor sabía algo nuevo sobre su propio libro: qué partes interesaban más, qué preguntas hacía el público, qué se vendía y qué no. La sabiduría era ancestral y simple: \textbf{el libro se completa con la audiencia, no antes}. La sustentación del periodo más la feria del libro escolar es esa práctica ancestral aplicada a tu grupo.
\end{softbox}

\begin{softbox}{milcTurquesa}{milcGris}{Saber contemporáneo}
La \textbf{sustentación pública} del libro es el acto formal donde el editor presenta su pieza ante el grupo. Tiene 4 propiedades irrenunciables: (1) \textbf{Tiempo controlado}: 6 minutos (similar a sustentación TED corta). Obliga a destilar lo esencial. (2) \textbf{Estructura clara}: concepción + proceso + decisiones + lectura de fragmento + muestra física o digital. (3) \textbf{Lectura en voz alta de un fragmento}: 1-2 minutos donde la audiencia escucha la voz del libro. Es la prueba de fuego: ¿suena a ti o suena a chatbot? (4) \textbf{Honestidad sobre uso de IA}: declarar abiertamente qué modelos usaste, qué porcentaje del texto generaste, qué porcentaje editaste o reescribiste a mano. La \textbf{feria del libro escolar} extiende la sustentación: los libros se exhiben, los compañeros se acercan a leer fragmentos, hay intercambio de comentarios. Es práctica de ciudadanía editorial: cada estudiante actúa como autor y como lector. La regla profesional del cierre: \textbf{el libro no termina cuando se firma; termina cuando se comparte con la audiencia}.
\end{softbox}

\begin{softbox}{milcMostaza}{milcGris}{Pregunta-puente}
\textit{¿Qué sabían los autores de la feria del libro del pueblo al leer fragmentos en voz alta a quien se acercaba a la mesa, que el estudiante novato olvida cuando entrega el PDF sin sustentar? ¿Y por qué la lectura en voz alta es la prueba de fuego de tu voz propia?}
\end{softbox}

\begin{softbox}{milcVerde}{milcGris}{Saber y saber hacer}
Al terminar podrás: (1) \textbf{analizar} tu libro completo identificando los 2-3 fragmentos más representativos de tu voz propia; (2) \textbf{explicar} en sustentación pública la concepción, proceso, decisiones críticas y uso de IA; (3) \textbf{evaluar} con honestidad las preguntas duras de la audiencia; (4) \textbf{crear} la presentación del libro y participar activamente en la feria como autor y como lector de los compañeros.
\end{softbox}



\small
\begin{tabularx}{\textwidth}{>{\bfseries}p{3.0cm}Y}
\rowcolor{milcGradePrimary!12}Autor & Dr. Álvaro Cárdenas Orozco\\
DBA & Producción editorial mediada por IA con criterio ético y técnico (MEN, T\&I)\\
Referentes & Ley 23 de 1982 · Floridi (ética de la IA generativa) · Dussel (autoría con dignidad)\\
Producto final & Sustentación pública de 6 minutos de tu libro frente al grupo con (a) presentación del libro (concepción + proceso + uso de IA + decisiones críticas), (b) lectura en voz alta de 1-2 minutos de un fragmento elegido, (c) muestra del PDF o libro impreso, (d) sesión de preguntas y respuestas de 5 minutos + participación activa en feria del libro escolar (puesto del libro, lectura de fragmentos de compañeros, intercambio honesto de comentarios) + autoevaluación final con 2 fortalezas y 2 mejoras\\
\end{tabularx}
\normalsize

\puente{Antes de empezar, mira el viaje completo. Has recorrido 9 sesiones para producir un libro de 80 páginas. Hoy es el \textbf{día de la firma pública}: presentas el libro ante el grupo. Después de hoy, cierra el periodo 1. El periodo 2 abre con informes técnicos y comunicación profesional.}

\milcestacion{milcGradeDark}{Ruta MILC: así vamos a aprender}
\begin{tabularx}{\textwidth}{>{\centering\arraybackslash}X>{\centering\arraybackslash}X>{\centering\arraybackslash}X>{\centering\arraybackslash}X}
\phasecard{1}{milcVerde}{milcGris}{Escucha\\[-1mm]\small Observo, escucho y pregunto.} &
\phasecard{2}{milcTurquesa}{milcGris}{Entender\\[-1mm]\small Sistematizo y ordeno.} &
\phasecard{3}{milcMagenta}{milcGris}{Hacer\\[-1mm]\small Construyo y aplico.} &
\phasecard{4}{milcVino}{milcGris}{Cierre\\[-1mm]\small Evalúo y reflexiono.}
\end{tabularx}


\puente{Antes de teorizar, vas a hacer una \emph{lectura crítica final} de tu libro completo: identifica 2-3 fragmentos donde sientas que tu voz propia se nota más. Esos son los candidatos a fragmento de lectura en sustentación. Léelos en voz alta para ti mismo: ¿suenan a ti?}

\milcestacion{milcVerde}{Estación 1 de 3 · Escucha}

\begin{softbox}{milcVerde}{milcGris}{Escena para escuchar}
\textbf{Actividad 1 · ANALIZA --- Selección de fragmentos para lectura en voz alta} (15 min · individual). Toma el PDF final de tu libro y elige 2-3 fragmentos de 1-2 minutos cada uno (aproximadamente 150-200 palabras) donde sientas que tu voz propia se nota más. Léelos en voz alta a ti mismo. Reglas: (1) eviten aperturas obvias de capítulo si la apertura fue muy IA. (2) busquen escenas con tu ejemplo personal, tu opinión clara, tu estilo. (3) lee cada candidato en voz alta y mide tiempo. Elige el que te enorgullezca más como editor responsable.
\end{softbox}

{\sffamily\bfseries\fontsize{8}{10}\selectfont\textcolor{milcNegro!55}{MARCA CUANDO LO HAYAS HECHO}}
\milccheckrow{Releí mi libro completo buscando fragmentos.}
\milccheckrow{Seleccioné 2-3 candidatos de 1-2 minutos.}
\milccheckrow{Leí cada uno en voz alta y elegí el final.}

\begin{infoband}{milcVerde}


\textbf{Lo que va al cuaderno (Actividad 1 · ANALIZA).} Título: «Actividad 1 --- Fragmentos para sustentar». \textbf{Formato:} 2-3 fragmentos identificados con página + tiempo de lectura + elección final con razón. \textbf{Sabes que terminaste cuando} el fragmento elegido te enorgullezca leerlo en público.
\end{infoband}

\puente{Ya tienes fragmentos candidatos. Ahora vas a entender la \textbf{anatomía de la sustentación de 6 minutos + feria}: estructura, los 5 errores típicos del sustentante novato, las reglas de la feria escolar.}

\milcestacion{milcTurquesa}{Estación 2 de 3 · Entender}

\begin{softbox}{milcTurquesa}{milcGris}{Lo que vas a entender}
Tu fragmento elegido es la prueba de tu voz propia. Ahora necesitas la \textbf{anatomía profesional} de la sustentación de 6 minutos y la feria escolar.
\end{softbox}

\milcpaso{1}{milcTurquesa}{La \textbf{estructura de la sustentación de 6 minutos} se descompone así: (1) \textbf{Concepción y producto} (1 min): qué libro hiciste, género, audiencia, propósito (ficha editorial de S2). (2) \textbf{Proceso y decisiones críticas} (1.5 min): cómo lo concebiste, cómo usaste prompts, cómo iteraste a V3, qué decisiones editoriales tomaste. (3) \textbf{Lectura en voz alta de fragmento} (1.5-2 min): el fragmento elegido en la actividad 1. (4) \textbf{Uso de IA declarado con honestidad} (1 min): qué modelos, qué porcentaje, qué hiciste tú. (5) \textbf{Muestra física o digital + cierre} (0.5-1 min): mostrar el PDF abierto, la portada, el lomo si es físico. Total: 6 minutos. Después: 5 minutos de Q\&A.}
\milcpaso{2}{milcTurquesa}{Los \textbf{5 errores típicos} del sustentante novato: (a) \textbf{Leer las slides}: matar la sustentación en el primer minuto. (b) \textbf{Ocultar uso de IA}: pretender escritura puramente humana cuando IA fue significativa. (c) \textbf{Fragmento débil}: leer un párrafo plano sin voz propia, traicionando todo el periodo. (d) \textbf{Evadir preguntas duras}: \emph{``¿qué porcentaje generó la IA?''} requiere respuesta honesta, no esquive. (e) \textbf{Cierre apurado}: terminar con \emph{``gracias''} en lugar de invitación a la feria.}
\milcpaso{3}{milcTurquesa}{Las \textbf{reglas de la feria del libro escolar} (después de las sustentaciones): (a) cada autor pone su libro en un puesto (mesa, escritorio). (b) los compañeros recorren los puestos en grupos de 3-4. (c) en cada puesto el autor lee 30 segundos de un fragmento y responde 1-2 preguntas. (d) los lectores pueden tomar nota de qué libros les interesaron más para leerlos completos después. (e) al cerrar la feria, cada autor recibe 3-5 comentarios escritos de compañeros sobre su libro.}
\milcpaso{4}{milcTurquesa}{Algoritmo: (1) preparar slides simples (5-7 slides máximo o usar el propio PDF como apoyo). (2) escribir guion hablado de 800-900 palabras (= 6 min). (3) ensayar 2 veces cronometrado. (4) sustentar. (5) leer fragmento en voz alta. (6) responder Q\&A con honestidad. (7) participar en feria como autor y como lector. (8) autoevaluar.}

\begin{drawbox}{milcTurquesa}{Anatomía de sustentación + feria}
\textbf{Sustentación (6 min):} concepción + proceso + fragmento + IA + muestra. \\[1mm] \textbf{Q\&A (5 min):} honestidad ante preguntas duras. \\[1mm] \textbf{Feria:} puesto + lectura 30 seg + 3-5 comentarios recibidos. \\[1mm] \textbf{Cierre:} autoevaluación con fortalezas y mejoras.
\end{drawbox}

\begin{softbox}{milcMostaza}{milcGris}{Errores comunes y su corrección}
(1) Leer las slides en vez de hablar. (2) Ocultar uso de IA. (3) Fragmento débil sin voz propia. (4) Evadir preguntas duras. (5) Cierre apurado sin invitar a la feria. (6) En la feria, no leer los libros de los compañeros (solo cuidar el propio puesto). (7) No tomar nota de los comentarios recibidos.
\end{softbox}

\begin{infoband}{milcTurquesa}


\textbf{Lo que va al cuaderno (Actividad 2 · EXPLICA).} Título: «Actividad 2 --- Anatomía sustentación + feria». \textbf{Formato:} (a) plan minuto a minuto de los 6 + 5 minutos; (b) reglas de la feria; (c) los 7 errores con marca del que más temes. \textbf{Sabes que terminaste cuando} tienes plan claro para mañana.
\end{infoband}

\puente{Tienes el método. Toca aplicarlo: preparas slides simples, ensayas cronometrado, sustentas en público, lees fragmento en voz alta, respondes preguntas, participas en feria, autoevalúas.}

\milcestacion{milcMagenta}{Estación 3 de 3 · Hacer}

\begin{softbox}{milcMagenta}{milcGris}{Cómo vas a construir tu producto}
\textbf{Actividad 3 · EVALÚA + CREA --- Sustentación ejecutada + feria + autoevaluación} (90 min total · individual + colectivo). Hora de ejecutar. Preparas slides, ensayas, sustentas con lectura de fragmento, respondes Q\&A, participas en feria como autor y como lector, autoevalúas.
\end{softbox}

\milcpaso{1}{milcMagenta}{Tu sustentación se construye en 6 pasos: (1) preparar slides simples (5-7 slides o usar el propio PDF como apoyo visual). (2) escribir guion hablado en prosa de 800-900 palabras. (3) ensayar cronometrado al menos 2 veces. (4) sustentar en público con lectura de fragmento elegido. (5) responder 3-5 preguntas con honestidad técnica. (6) participar en feria como autor (puesto con tu libro) y como lector (recorrer los demás).}
\milcpaso{2}{milcMagenta}{Sigue el patrón \textbf{``honestidad pública sobre uso de IA''}: la declaración del uso de IA en sustentación es momento delicado pero esencial. Modelo de respuesta: \emph{``Usé Bing Image Creator para la portada, ChatGPT para borradores de los capítulos 2-5, y escribí enteramente a mano los capítulos 1, 6 y la carta del editor. Aproximadamente el 40 por ciento del texto final pasó por IA, el 60 por ciento fue escrito o reescrito por mí''}. Esa precisión gana respeto.}
\milcpaso{3}{milcMagenta}{La \textbf{autoevaluación final} tiene 4 elementos: (a) \textbf{2 fortalezas}: qué hiciste bien en la sustentación (manejo del tiempo, fragmento elegido, claridad de la declaración de IA). (b) \textbf{2 mejoras}: qué harías distinto en una próxima sustentación (más ensayo, mejor manejo del nervio, slides más simples). (c) \textbf{1 cosa aprendida del feedback de la feria}: qué te dijeron los compañeros que no esperabas. (d) \textbf{Compromiso para el periodo 2}: 1 acción concreta para el siguiente trimestre.}
\milcpaso{4}{milcMagenta}{Algoritmo: (1) preparar slides. (2) ensayar. (3) sustentar. (4) responder Q\&A. (5) participar en feria 30 min (15 min como autor, 15 min como lector). (6) recoger comentarios recibidos. (7) autoevaluar al cerrar la sesión.}

\begin{drawbox}{milcMagenta}{Checklist antes de sustentar}
\checkbox Slides simples (5-7 o usando el PDF como apoyo). \\[1mm] \checkbox Guion hablado de 800-900 palabras. \\[1mm] \checkbox Fragmento elegido y cronometrado (1-2 min). \\[1mm] \checkbox 2 ensayos cronometrados hechos. \\[1mm] \checkbox Declaración de uso de IA preparada con porcentaje. \\[1mm] \checkbox Libro disponible para mostrar (PDF abierto o impreso). \\[1mm] \checkbox Puesto listo para la feria. \\[1mm] \checkbox Autoevaluación pendiente al cerrar.
\end{drawbox}

\begin{softbox}{milcMagenta}{milcGris}{Plantilla de trabajo}
\textbf{Concepción (1 min).} \textit{Género + tema + audiencia + propósito.} \\[1mm] \textbf{Proceso (1.5 min).} \textit{Concepción + prompts + iteración + decisiones.} \\[1mm] \textbf{Lectura de fragmento (1.5-2 min).} \textit{Voz propia en escena.} \\[1mm] \textbf{Uso de IA (1 min).} \textit{Modelos + porcentajes + reparto humano.} \\[1mm] \textbf{Muestra y cierre (0.5-1 min).} \textit{PDF + invitación a la feria.} \\[1mm] \textbf{Q\&A (5 min).} \textit{Honestidad ante preguntas duras.}
\end{softbox}

\begin{infoband}{milcMagenta}


\textbf{Lo que va al cuaderno (Actividad 3 · EVALÚA + CREA).} Título: «Actividad 3 --- Mi sustentación + feria». \textbf{Formato:} (a) referencia a slides o PDF apoyo; (b) guion hablado; (c) preguntas recibidas con respuesta; (d) comentarios de la feria; (e) autoevaluación con 4 elementos. \textbf{Sabes que terminaste cuando} entregaste libro al docente y reservaste tiempo para reflexión.
\end{infoband}

\puente{Lo que sigue resume \emph{qué} entregas hoy: sustentación ejecutada + feria participada + autoevaluación. El criterio principal: que la audiencia se vaya sabiendo qué libro hiciste y queriendo abrirlo.}

\milcestacion{milcMostaza}{Producto de la sesión}
\textbf{Sustentación 6 min + lectura de fragmento + feria participada + autoevaluación}.
\begin{softbox}{milcMostaza}{milcGris}{Criterios de calidad}
(1) Sustentación dentro de 6 minutos. (2) Estructura clara de 5 segmentos. (3) Lectura de fragmento en voz alta. (4) Declaración honesta de uso de IA con porcentajes. (5) Participación en feria como autor y lector. (6) Autoevaluación con 4 elementos.
\end{softbox}

\puente{Antes de cerrar, mira la sustentación desde las cinco dimensiones humanas. Defender el libro con honestidad y orgullo es cierre del oficio editorial; entregar sin sustentar es desperdicio del periodo entero.}

\milcestacion{milcVino}{Evaluación liberadora · cinco dimensiones}

\begin{softbox}{milcVino}{milcGris}{Sentido de la evaluación}
Evaluar no es solo revisar una nota. Es reconocer qué aprendí, qué cambió en mí, qué emociones manejé, cómo me relaciono con la ciudad y la comunidad, y qué le dejaré a quien viene detrás.
\end{softbox}

\textbf{Mi reflexión liberadora en cinco dimensiones.} Completa cada frase iniciada:

\small
\begin{tabularx}{\textwidth}{>{\bfseries\RaggedRight\arraybackslash}p{3.4cm}Y>{\RaggedRight\arraybackslash}p{4.6cm}}
\rowcolor{milcVino}\color{white}Dimensión & \color{white}\bfseries Mi respuesta (frame starter) & \color{white}\bfseries Algo que descubrí\\
1. Personal & Lo que más cambió en mí fue \linead{6.5cm} & \linead{4.2cm}\\
2. Emocional & La emoción más fuerte fue \linead{2.5cm} y la regulé \linead{3cm} & \linead{4.2cm}\\
3. Ciudadana & En democracia esto importa porque \linead{6cm} & \linead{4.2cm}\\
4. Local (Cartago) & En mi barrio sería útil para \linead{6.5cm} & \linead{4.2cm}\\
5. Intergeneracional & A mi abuela/abuelo le diría \linead{6.5cm} & \linead{4.2cm}\\
\end{tabularx}
\normalsize

\begin{infoband}{milcVino}
\textbf{Para la próxima sustentación.} Vuelve a esta tabla en seis meses. Verás que las respuestas cambian — no porque lo aprendido cambie, sino porque tú cambias. La evaluación liberadora es una conversación contigo a través del tiempo.
\end{infoband}


\puente{Y para terminar, tres voces sobre el cierre del oficio editorial. Dussel sobre el editor que devuelve a la comunidad lo aprendido, Marco Aurelio sobre la honestidad del cierre, Floridi sobre la sustentación como ética profesional contemporánea.}

\milcestacion{milcGradeDark}{Triángulo de pensamiento · cierre filosófico}

\begin{softbox}{milcVino}{milcGris}{Dussel · lente del nosotros}
\textit{El editor que sustenta su libro ante la comunidad devuelve oficio asumido, no archiva trabajo individual.}

{\footnotesize\itshape\color{milcNegro!70}--- formulación de esta guía, al modo de Enrique Dussel. No es una cita textual.}

Tu sustentación cierra el ciclo del periodo: el libro vuelve a la comunidad (el grupo) hecho voz propia firmada. La filosofía de la liberación pide ese cierre: \emph{el saber editorial no se guarda; se devuelve}. La feria del libro es ese acto de devolución colectiva. Cuando lees tu fragmento y muestras tu libro, estás haciendo lo que el editor de la gaceta barrial hacía cada semana: ofrecer la pieza terminada al público que la espera.

\textbf{¿Mi sustentación devuelve oficio asumido al grupo, o solo cumple con requisito escolar?}
\end{softbox}
\linead{\linewidth}\\[1mm]
\linead{\linewidth}

\begin{softbox}{milcMostaza}{milcGris}{Estoicismo · Marco Aurelio · cuidado interior}
\textit{La honestidad del cierre es virtud del editor; ocultar el proceso es debilidad disfrazada de modestia.}

{\footnotesize\itshape\color{milcNegro!70}--- formulación de esta guía, al modo de Marco Aurelio. No es una cita textual.}

Es tentador minimizar el uso de IA al sustentar (\emph{``casi no usé IA''}) o exagerarlo (\emph{``todo lo hizo la IA''}). La sabiduría estoica recomienda lo opuesto: \emph{declara los porcentajes exactos con sobriedad}. Esa precisión gana respeto profesional que las versiones inflada o minimizada nunca ganan. La phronesis del cierre honra la honestidad numérica.

\textbf{¿Estoy declarando porcentajes honestos del uso de IA, o inflando o minimizando según convenga?}
\end{softbox}
\linead{\linewidth}\\[1mm]
\linead{\linewidth}

\begin{softbox}{milcTurquesa}{milcGris}{Floridi · lente de la infoesfera}
\textit{La sustentación con declaración honesta de IA es la nueva ética profesional del oficio editorial del siglo XXI.}

{\footnotesize\itshape\color{milcNegro!70}--- formulación de esta guía, al modo de Luciano Floridi. No es una cita textual.}

En la era de la IA generativa, los editores profesionales que declaran abiertamente el uso de IA en sus libros se posicionan en el modelo ético emergente. Quien oculta uso de IA acumula deuda ética que tarde o temprano se descubre. Tu sustentación de hoy te entrena en una práctica que rige toda producción editorial profesional emergente: \textbf{transparencia con porcentajes y modelos}.

\textbf{¿Mi sustentación contribuye al modelo editorial transparente del siglo XXI, o sostiene la opacidad heredada?}
\end{softbox}
\linead{\linewidth}\\[1mm]
\linead{\linewidth}

\begin{infoband}{milcNegro}
\textbf{Nota para el docente.} Las tres formulaciones son del autor de la guía, no citas textuales de los pensadores: recogen su lente para pensar el tema, no sus palabras. Si vas a citarlos en clase, remite a la obra original.
\end{infoband}

\milcestacion{milcVino}{Mi autoevaluación · marca una casilla por fila}

{\small
\setlength{\extrarowheight}{2.2pt}
\begin{tabularx}{\textwidth}{>{\RaggedRight\arraybackslash\bfseries}p{3.0cm}YYY}
\rowcolor{milcVino}\color{white}\bfseries Criterio & \color{white}\bfseries Lo logré & \color{white}\bfseries En proceso & \color{white}\bfseries Necesito apoyo\\[.6mm]
Comprensión del tema & \milcopcion{Explico las ideas con mis palabras.} & \milcopcion{Reconozco las ideas, falta precisión.} & \milcopcion{Necesito repasar lo básico.}\\[.8mm]
\rowcolor{milcGris}Pensamiento computacional & \milcopcion{Usé los pilares al armar mi producto.} & \milcopcion{Usé algunos pilares.} & \milcopcion{Necesito releer la estación 2.}\\[.8mm]
Aplicación y producto & \milcopcion{Mi producto cumple los criterios.} & \milcopcion{Está armado, con ajustes pendientes.} & \milcopcion{Aún está en construcción.}\\[.8mm]
\rowcolor{milcGris}Reflexión 5 dimensiones & \milcopcion{Respondí cada una con honestidad.} & \milcopcion{Respondí algunas con detalle.} & \milcopcion{Mis respuestas son muy generales.}\\[.8mm]
Triángulo de pensamiento & \milcopcion{Conecté las 3 voces con mi trabajo.} & \milcopcion{Conecté al menos una.} & \milcopcion{No las relaciono con mi proceso.}\\[.8mm]
\end{tabularx}}

\vspace{3mm}
\textbf{Hoy entendí que:}\\[1mm]
\linead{\linewidth}\\[1mm]
\linead{\linewidth}

\textbf{Mi compromiso para la próxima sesión es:}\\[1mm]
\linead{\linewidth}\\[1mm]
\linead{\linewidth}

\begin{softbox}{milcMostaza}{milcGris}{Cita de despedida · Marco Aurelio}
\textit{``El obstáculo en el camino se vuelve el camino. Lo que estorba al camino se vuelve camino.''}\\[1mm]
Cada error de hoy, bien observado, se convierte en pieza del camino que pisarás mañana.
\end{softbox}

\small
\begin{tabularx}{\textwidth}{>{\bfseries}p{3.6cm}Y}
\rowcolor{milcGradeDark!12}Nombre completo: & \linead{10cm}\\
Fecha: & \linead{4cm} \quad Curso/grupo: \linead{4cm}\\
Mi firma: & \linead{9cm}\\
\end{tabularx}
\normalsize

\begin{infoband}{milcGradeDark}
\textbf{Para la próxima sesión.} Cierras el periodo 1 (escritura abierta con IA) y abre el periodo 2 (informes técnicos y comunicación profesional con IA). La sesión 1 del periodo 2 trabajará los \textbf{4 tipos de texto profesional}: del libro creativo a los informes que se usan en empresas, ONGs y administración pública.
\end{infoband}

\end{document}
